\section{Derivation of equations and results}
\label{sec:derivations}

The closed forms of Section~\ref{sec:analytical} all descend from one Riccati
equation, obtained twice over --- once for $I$ and once for $D$ --- with
different initial conditions. We derive that equation, solve it, and then
assemble the remaining results from it.

\subsection{The Riccati equation for $I(t)$}

Considering what may happen during an interval of duration $\delt$ after
$t=0$, we can write down
$$I(0+\delt) = \mu \delt + [1 - (\lambda + \mu + \delta)\delt] I(0) + \lambda\delt\, {I(0)}^2 + o(\delt).$$
The single initial unit may die, divide, cause rupture, or do nothing at all;
each possibility is accounted for on the right-hand side, with the respective
event probabilities multiplied by the updated non-rupture likelihoods --- as if
the process were beginning anew. In the case that a rupture occurred, this value
is obviously $0$; if the unit dies, non-rupture is certain, contributing the
$\mu\delt$ term; if a birth takes place, the non-rupture probability assumes
the value $I_2(0)$, which, by the branching property, equals ${I(0)}^2$.
Because the model is homogeneous in time, the relation maintains applicability
in subsequent intervals $[t, t+\delt]$:
$$I(t+\delt) = \mu \delt + [1 - (\lambda + \mu + \delta)\delt] I(t) + \lambda\delt\, {I(t)}^2 + o(\delt).$$
Rearranging yields
$$\frac{I(t+\delt) -I(t)}{\delt} = \mu - (\lambda + \mu + \delta) I(t) + \lambda {I(t)}^2 + \frac{o(\delt)}{\delt},$$
which, in the limit as $\delt\to 0$, becomes
\begin{equation}
\label{ddII}
\ddt{I} = \mu - (\lambda + \mu + \delta) I + \lambda {I}^2,
\end{equation}
a Riccati equation. Its quadratic right-hand side factorises through the roots
$a$ and $b$ defined in Section~\ref{sec:reference}, giving the equivalent form
\begin{equation}
\label{ddIIfact}
\ddt{I} =\lambda\lrb{I-a}\lrb{I-b}.
\end{equation}
Plotting $\dd I/\dd t$ against $I$ gives an intuitive picture of why $I(t)$
tends to $b$ at late time: $b$ is the stable root, and any trajectory starting
at $I(0)=1<a$ runs downhill towards it.

The identities collected in \Eqref{eq:identities} fall out of this
factorisation. Expanding $\lambda(I-a)(I-b)$ and comparing coefficients with
\eqref{ddII} gives, by Vieta's formulas,
$$ab=\frac{\mu}{\lambda},\qquad a+b=\frac{\lambda+\mu+\delta}{\lambda}.$$
The two derived constants $A=a-1$ and $B=1-b$ then combine as
$$AB = (a-1)(1-b) = a + b - ab - 1
= \frac{\lambda+\mu+\delta}{\lambda} - \frac{\mu}{\lambda} - 1
= \frac{\delta}{\lambda},$$
while $A+B=(a-1)+(1-b)=a-b$ and $\theta=\lambda(a-b)$ are immediate from the
definitions. These four identities are the bookkeeping that makes every closed
form in the chapter collapse to its shortest shape, and they are worth
committing to memory.

\subsection{Solving the Riccati equation}

Since the right-hand side of \eqref{ddIIfact} is a product of linear factors,
the equation separates:
$$\int \frac{\dd I}{(I-a)(I-b)} = \lambda t + \text{const}.$$
The partial-fraction decomposition
$$\frac{1}{(I-a)(I-b)} = \frac{1}{a-b}\lrb{\frac{1}{I-a}-\frac{1}{I-b}}$$
integrates the left-hand side to $(a-b)^{-1}\log\frac{I-a}{I-b}$, so every
solution satisfies
\begin{equation}
\label{eq:riccati_ratio}
\frac{I(t)-a}{I(t)-b} = \frac{I(0)-a}{I(0)-b}\,\ee^{\lambda(a-b)t} = \frac{I(0)-a}{I(0)-b}\,w.
\end{equation}
With $I(0)=1$, the ratio on the right is $(1-a)/(1-b)=-A/B$, and solving
\eqref{eq:riccati_ratio} for $I(t)$ reproduces the closed form of
Theorem~\ref{thm:I}:
\begin{equation}
\label{formI}
I(t) = \frac{aB + bA\,w}{B + A\,w}.
\end{equation}
(Some treatments of this process label the two roots $I_+$ and $I_-$; the
correspondence is $I_+=a$ and $I_-=b$, and we use $a$, $b$ throughout.)

\subsection{The internal-extinction probability $D(t)$}

The probability $D(t)=p_0(t)$ that the process has died out internally by time
$t$ satisfies the \textit{same} Riccati equation as $I$, with a different
initial condition. The first-step decomposition is identical in structure: if
the first unit dies, the process sits at $0$ with certainty; if it divides,
both descendant families must die out, an event of probability $D(t)^2$; if it
causes a rupture, the process is at $H$ and not at $0$. Hence
\begin{equation}
\label{ddD}
\ddt{D} = \mu - (\lambda + \mu + \delta) D + \lambda D^2,\qquad D(0)=0.
\end{equation}
In the ratio form \eqref{eq:riccati_ratio}, the initial ratio is now
$(0-a)/(0-b)=a/b$, so
$$\frac{D(t)-a}{D(t)-b} = \frac{a}{b}\,w,$$
and solving for $D(t)$ gives the closed form of Theorem~\ref{thm:D}:
$$D(t) = \frac{ab(w-1)}{aw-b}.$$
That $I$ and $D$ obey the same differential equation is not a coincidence: both
are ``survival-type'' quantities for the same process, distinguished only by
which absorbing state is counted as survival, and the shared equation is what
makes their difference simplify so cleanly.

\subsection{The non-fixation probability $\Ifix(t)$}

With $I$ and $D$ in hand, $\Ifix=I-D$ could be written down immediately; but the
direct derivation is worth recording, because it exposes the structure that the
multi-founder formulas of Section~\ref{sec:analytical} inherit. A first-step
decomposition for $\Ifix$ differs from the two above in the birth branch: after
a division, non-fixation of the pair requires that \textit{neither} family
fixates, whether by death or by rupture, and that event has probability
$I^2-D^2$, not $\Ifix^2$. Thus
\begin{equation}
\label{ddIfix}
\ddt{\Ifix} = -(\lambda + \mu + \delta)\Ifix + \lambda\lrb{I^2 - D^2},\qquad \Ifix(0)=1,
\end{equation}
where the death and catastrophe branches contribute nothing, each sending the
process to one of the absorbing states. Subtracting \eqref{ddD} from
\eqref{ddII} yields the same equation, confirming $\Ifix = I - D$. The closed
form follows by algebra:
$$\Ifix(t) = I(t) - D(t) = \frac{aB + bAw}{B+Aw} - \frac{ab(w-1)}{aw-b} = \frac{(a-b)^2\,w}{(B+Aw)(aw-b)}.$$
The limits read off from this expression agree with the definition:
$\Ifix(0)=(a-b)^2/[(A+B)(a-b)]=1$, and $\Ifix(t)\to0$ as $w\to\infty$, since
the numerator grows like $w$ while the denominator grows like $w^2$.
Differentiating \eqref{ddIfix} at $t=0$ gives the initial slope
$\Ifix'(0)=-(\lambda+\mu+\delta)+\lambda=-(\mu+\delta)$, the total rate at
which a productive realisation is lost at the first instant. When $\mu=0$ we
have $b=0$ and $D\equiv0$, and \eqref{ddIfix} collapses to the Riccati
equation for $I$ itself.

\subsection{The forward relation $I'=-\delta J$}
\label{sec:forwardI}

One further relation, connecting the non-catastrophe probability to the mean
load, will be used repeatedly. The forward Kolmogorov equation for the
probability of the rupture state is
$$\ddt{}\pr{\xt = H} = \delta\sum_{n\ge1} n\,p_n(t) = \delta J(t):$$
the state $H$ is entered from every productive state $n$, at rate $\delta n$.
Since $I(t)=1-\pr{\xt=H}$, this gives immediately
\begin{equation}
\label{ddIJ}
\ddt{I} = -\delta J,
\end{equation}
which is the first of the key ODEs of Section~\ref{sec:reference} in a second
guise, and from which $J(t)=-\delta^{-1}I'(t)$ follows by differentiation of
\eqref{formI}, as claimed in Theorem~\ref{thm:J}.
